\section{Introduction}\label{sec:intro}

Persistent homology has two software cultures. One writes fast specialists. Ripser~\cite{bauer2021ripser} computes
Vietoris--Rips barcodes with a small memory footprint; Cubical Ripser~\cite{kaji2020cubical}, Flash
Cubical~\cite{lebreton2026flash} and RedZeD~\cite{kapulkin2026redzed} each take one kind of input, mostly over $\Ftwo$, and
report bars. The other writes toolkits: GUDHI~\cite{gudhi2014}, Dionysus~\cite{dionysus2}, Eirene~\cite{henselman2016matroid},
Ripserer.jl~\cite{cufar2020ripserer}, OAT~\cite{hang2021umatch}, giotto-tda~\cite{tauzin2020gtda} and JavaPlex~\cite{tausz2014javaplex},
which cover more complexes and more of the pipeline around a barcode. The two cultures answer different questions. A
specialist answers ``how fast can this barcode be computed?'' A toolkit answers ``which of the many ways to build a
filtered complex, reduce it, and read the result will let me finish my analysis?''

\tdaj{} is a toolkit that takes the specialists' techniques seriously. It began as the successor to JavaPlex, the
research package for persistent (co)homology that grew out of the Computational Topology workgroup at Stanford and was
used mostly through MATLAB~\cite{tausz2014javaplex}. More than a decade later the needs are different: users call persistent homology
from Python and from the browser, expect a barcode in seconds for the point-cloud sizes that Ripser made routine, and
increasingly want to see \emph{where} a class lives, with a cycle or cocycle they can plot, project or feed into
circular coordinates, rather than only that it exists. \tdaj{} is built for those users, and as a platform on which to try
new ideas on persistent homology itself: the machinery behind each entry point is public, so a new filtration, a new
reduction, or a new pairing strategy can be tried against the library's oracles without forking it.

\paragraph{What the library is.}
The core is a Scala~3 library for the JVM and for JavaScript, built from one source tree. A filtered complex is a stream
of cells in filtration order; an engine consumes the stream and returns a barcode with representatives; a verb,
\code{Persistence(input, \ldots)}, picks sensible defaults for both. Everything between input and barcode is generic over the
cell type and over the coefficient field, and the field is a type, so mixing elements of two different prime fields is a compile-time
error rather than a silently wrong answer (\cref{sec:design}). Beyond barcodes the library offers circular and toroidal
coordinates, zigzag and relative persistence, maps between barcodes, discrete Morse reduction, the persistent Laplacian,
barcode distances and vectorizations, and simplicial sets with products, cup products and Steenrod squares
(\cref{sec:breadth}).

\paragraph{Contributions.}
This is a system paper. We separate what is borrowed from what is new, and
we say where a claim waits on a measurement.
\begin{enumerate}[label=\textbf{C\arabic*.}, leftmargin=2.2em, itemsep=2pt]
\item \textbf{A typed substrate for persistent homology} (\cref{sec:design}): cells, chains, fields, filtered complexes and
  engines as separate, checked abstractions; a single verb on top; and defaults that encode what went wrong in practice
  ($\Fsev$ rather than $\Ftwo$, zero-length bars dropped, the enclosing radius as the threshold).
\item \textbf{A representative for every bar, from every engine} (\cref{sec:representatives}): each engine computes cycles or
  cocycles natively and derives the other kind from its own pairing by one reduction pass, in both directions, with every pivot
  checked against the pairing. The derivation of cycles from cohomology is due to \v{C}ufar and
  Virk~\cite{cufar2021involuted}; we make it uniform over cell types and fields, add the reverse direction, and make
  representatives lazy so that they are paid for only when read.
\item \textbf{An engine family built from a few techniques} (\cref{sec:algorithms}): in-order generation of
  Vietoris--Rips simplices from the preprint of Bauer, Matuszewski and Vejdemo-Johansson~\cite{bauer2024generating}; a packed
  Ripser engine with one- or two-word simplex indices; degree $0$ by union-find in every engine that reduces it; a dual
  union-find engine for the top degree of cubical and alpha complexes that extends the idea of Flash
  Cubical~\cite{lebreton2026flash} to the $T$-construction, to any field and to Delaunay complexes, with signed
  representatives; and exact parallel reduction in units of edge blocks.
  \PROV{Mikael: which parts of the edge-block material are in the preprint and which are new here; see Appendix~\ref{app:notes}, question~1.}
\item \textbf{A correctness methodology and a cross-library agreement check} (\cref{sec:correctness}): reference engines kept in the
  test tree and compared term for term, oracles chosen so that the plausible wrong answer fails, and a benchmark harness that checks
  every tool's barcode against the others before it times them. The harness found a bug in our own default alpha backend and
  defects in published algorithms and reference code.
\item \textbf{One library, six front ends} (\cref{sec:platforms}): the JVM and JavaScript builds share their sources;
  MATLAB, Java, Python, the command line and npm are thin layers over one option set; and every number in the
  tutorials is asserted by a test.
\item \textbf{A same-machine evaluation against \TODOnum{number of tools} tools} (\cref{sec:evaluation}), reported as measured,
  including where specialists are faster.
\end{enumerate}

\paragraph{Summary of the evaluation.}
\OUTCOME{On Vietoris--Rips complexes, the default engine returning cycles is \TODOnum{factor range} slower than
ripser.cpp, and returns a representative for every bar; the packed engine computing cocycles alone is within
\TODOnum{factor range}. Parallel reduction in edge blocks gives \TODOnum{factor range} on \TODOnum{number} threads.
On cubical images the dual union-find engine is \TODOnum{factor range} from Cubical Ripser. On alpha complexes the
triangulation backends and the dual-quadratic-programming backend cross over at ambient dimension \TODOnum{d}.
The Scala.js build is \TODOnum{factor range} slower than the JVM build.}

\paragraph{Outline.}
\Cref{sec:background} fixes notation and places \tdaj{} among the existing software.
\Cref{sec:design} describes the design and \cref{sec:algorithms} the algorithms and the engineering behind the engine family.
\Cref{sec:correctness} describes how correctness is established, and \cref{sec:evaluation} the evaluation.
\Cref{sec:showcases} gives worked examples, \cref{sec:limitations} lists limitations, and \cref{sec:availability}
says where to find the software and how AI assistance was used in writing the code and this text.
